\subsection{Branching trees send piecewise-affine Schur responses}
\label{subsec:tree-response}

The path proof uses one realized scalar message at each depth.  On a
nonsymmetric tree, a subtree must describe how its root reacts as the parent
value changes.  The next theorem shows that this response remains highly
structured.

Root a tree at the single seed $o$.  For a nonroot vertex $u$ with parent
$p$, let $T_u$ be the descendant subtree rooted at $u$ and put
$\beta_{pu}:=-Q_{pu}>0$.  Given a hypothetical parent value $t\geq0$, define
the conditional subtree solution
\begin{equation}
 z^{(u)}(t)
 :=\argmin_{z\in\mathbb R_+^{T_u}}
 \left\{
  \frac12 z^\top Q_{T_uT_u}z-c_{\rho,T_u}^\top z
  -\beta_{pu}t z_u
 \right\},
 \qquad
 \mu_u(t):=\beta_{pu}z_u^{(u)}(t).
 \label{eq:conditional-subtree-response}
\end{equation}
Here $\mu_u(t)$ is the nonnegative Schur response returned to the parent.

\begin{theorem}[Monotone piecewise-affine tree response]
\label{thm:tree-response}
For every nonroot subtree $T_u$:
\begin{enumerate}[label=(\roman*),leftmargin=*]
 \item $z^{(u)}(t)$ is continuous and coordinatewise nondecreasing in $t$;
 its supports are nested.
 \item With the explicit activation threshold
 \begin{equation}
  \tau_u
  :=\frac{\alpha\rho\sqrt{d_u}}{\beta_{pu}},
  \label{eq:subtree-activation-threshold}
 \end{equation}
 one has $z^{(u)}(t)=0$ if and only if $0\leq t\leq\tau_u$.
 \item The response $\mu_u$ is continuous, nondecreasing, and piecewise
 affine with at most $|T_u|+1$ pieces.  On an interval with fixed nonempty
 active set $A\subseteq T_u$,
 \begin{align}
  z_A^{(u)}(t)
  &=Q_{AA}^{-1}
    \bigl(c_{\rho,A}+\beta_{pu}t e_u\bigr),
  \label{eq:tree-affine-piece}\\
  \mu_u'(t)
  &=\beta_{pu}^2(Q_{AA}^{-1})_{uu}>0.
  \label{eq:tree-response-slope}
 \end{align}
 \item If $c_{\rho,o}>0$, the root value of the full RPPR solution is the
 unique positive solution of
 \begin{equation}
  Q_{oo}t-c_{\rho,o}
  -\sum_{u:\,p(u)=o}\mu_u(t)=0.
  \label{eq:root-response-equation}
 \end{equation}
 On every fixed-support interval, the left-hand side has positive slope equal
 to the Schur complement of the active descendants in $Q$.  In particular,
 \begin{equation}
  \alpha
  \leq Q_{oo}-\sum_{u:\,p(u)=o}\mu_u'(t)
  \leq Q_{oo},
  \label{eq:root-response-slope-bound}
 \end{equation}
 wherever the derivatives exist.
\end{enumerate}
\end{theorem}

\begin{proof}
Strong convexity gives a unique conditional solution and continuity in the
linear parameter $t$.  If $t_2\geq t_1$, then the demand in the $u$ coordinate
increases by $\beta_{pu}(t_2-t_1)$.  The obstacle comparison principle for a
Stieltjes matrix gives
$z^{(u)}(t_2)\geq z^{(u)}(t_1)$, proving monotonicity and nested support.

Every vertex of $T_u$ is a nonseed, so
$c_{\rho,v}=-\alpha\rho\sqrt{d_v}<0$.  At the zero point, every descendant
$v\neq u$ has strictly positive obstacle gradient
$-c_{\rho,v}$, while the $u$ gradient is
$-c_{\rho,u}-\beta_{pu}t
=\alpha\rho\sqrt{d_u}-\beta_{pu}t$.
Thus zero satisfies all conditional KKT inequalities exactly when
$t\leq\tau_u$.  When $t>\tau_u$, the $u$ inequality is violated, so the
solution is nonzero.  Moreover, every nonempty active component must contain
$u$: a component separated from $u$ has only the strictly negative shifted
demand $c_{\rho,C}$ and would satisfy
$Q_{CC}z_C=c_{\rho,C}<0$, contradicting
$Q_{CC}^{-1}\geq0$ and $z_C>0$.  Hence $u$ itself is active whenever the
conditional solution is nonzero.  This proves (ii).

On an interval where the active set is $A$, its positive coordinates solve
the linear system in \cref{eq:tree-affine-piece}.  Principal Stieltjes
inverses are nonnegative, giving \cref{eq:tree-response-slope}.  By (i), an
active vertex never leaves as $t$ increases; every change of affine formula
therefore activates at least one previously inactive vertex.  There are at
most $|T_u|$ such events, proving (iii).

Finally, conditioned on root value $t$, the child subtrees are independent
and return $\mu_u(t)=\beta_{ou}z_u^{(u)}(t)$.  The active root KKT equality is
exactly \cref{eq:root-response-equation}.  On a fixed global active set,
eliminating all active descendants leaves the scalar Schur complement at the
root; it is positive because $Q$ is positive definite.  Hence the left-hand
side is strictly increasing piece by piece and its zero is unique.  Because
the relevant principal block is bounded below by $\alpha I$, the scalar Schur
complement is at least $\alpha$, proving
\cref{eq:root-response-slope-bound}.
\end{proof}

\begin{corollary}[Exact inverse-response recursion]
\label{cor:inverse-tree-response}
For a nonroot vertex $u$, let $\mathcal C(u)$ be its children and define
\begin{equation}
 H_u(z)
 :=\frac{
  Q_{uu}z-c_{\rho,u}
  -\sum_{v\in\mathcal C(u)}\mu_v(z)
 }{\beta_{pu}},
 \qquad z\geq0.
 \label{eq:inverse-response-map}
\end{equation}
Then $H_u(0)=\tau_u$, $H_u$ is continuous piecewise affine, and on every
linear piece
\begin{equation}
 H_u'(z)
 =\frac{Q_{uu}-\sum_{v\in\mathcal C(u)}\mu_v'(z)}{\beta_{pu}}
 \geq\frac{\alpha}{\beta_{pu}}>0.
 \label{eq:inverse-response-slope}
\end{equation}
Consequently,
\begin{equation}
 z_u^{(u)}(t)
 =\begin{cases}
  0,&0\leq t\leq\tau_u,\\
  H_u^{-1}(t),&t>\tau_u,
 \end{cases}
 \qquad
 \mu_u(t)=\beta_{pu}z_u^{(u)}(t).
 \label{eq:tree-response-inverse}
\end{equation}
The breakpoints used to construct $H_u$ are exactly the union of the child
response breakpoints, plus the activation point at zero; branching introduces
merging but no combinatorial proliferation.
\end{corollary}

\begin{proof}
Condition on $z_u=z$.  Each child subtree minimizes independently and sends
back $\mu_v(z)=\beta_{uv}z_v^{(v)}(z)$.  The active $u$ equation is
\[
 Q_{uu}z-c_{\rho,u}-\beta_{pu}t
 -\sum_{v\in\mathcal C(u)}\mu_v(z)=0,
\]
which is equivalent to $t=H_u(z)$.  On a fixed collection of active child
pieces, the numerator slope is the Schur complement obtained by eliminating
their active descendants.  The same argument as
\cref{eq:root-response-slope-bound} bounds it below by $\alpha$.  Thus $H_u$
is strictly increasing and invertible above $H_u(0)=\tau_u$, proving the
recursion.  Addition changes slopes only at child breakpoints, and inversion
maps those breakpoints monotonically without creating new ones.
\end{proof}

\begin{corollary}[Polynomial exact solver on an arbitrary finite tree]
\label{cor:finite-tree-message-solver}
Let $n=|V|$, let $\Delta$ be the maximum degree, and root a finite tree at the
single seed.  Explicitly materializing every subtree response from the leaves
upward, solving \cref{eq:root-response-equation}, and evaluating the responses
top down returns the exact RPPR solution in
\begin{equation}
 O\!\left(
  \log(1+\Delta)\sum_{u\neq o}|T_u|
 \right)
 \subseteq O(n^2\log(1+\Delta))
 \label{eq:finite-tree-message-work}
\end{equation}
arithmetic, with $O(\sum_{u\neq o}|T_u|)\subseteq O(n^2)$ response storage.
There is no exponential active-set enumeration.
\end{corollary}

\begin{proof}
A leaf response has one inactive and one affine piece.  Inductively, suppose
the sorted response breakpoints of all children of $u$ are available.  Their
arguments are already the common parent value $z_u$, so a multiway merge
forms the sorted breakpoint list of the sum in
\cref{eq:inverse-response-map}.  On every resulting interval, update the
aggregate slope and intercept of $H_u$ and invert that affine segment.  By
\cref{eq:inverse-response-slope}, the ordering is preserved.  The child lists
contain at most
$\sum_{v\in\mathcal C(u)}|T_v|=|T_u|-1$ breakpoints, so the work at $u$ is
$O(|T_u|\log(1+d_u))$ and the stored output has $O(|T_u|)$ pieces.  Sum over
vertices.  The root uses the same merge to locate the unique zero of
\cref{eq:root-response-equation}; one top-down evaluation then recovers every
coordinate.
\end{proof}

The preceding corollary materializes response pieces that lie far beyond the
eventual root value.  Locality instead requests child events lazily and stops
as soon as the root equation is solved.  The price of the elementary lazy
implementation is one propagation through the active ancestors per realized
event.  The next lemma shows that this depth factor is already on the
$1/\sqrt\alpha$ scale.

\begin{lemma}[Graph-universal Chebyshev support-radius bound]
\label{lem:graph-support-radius}
Let
\begin{equation*}
 R_\star:=\max\{\operatorname{dist}(o,v):v\in S^\star(\rho)\},
\end{equation*}
with $R_\star=0$ when the support is empty or contains only the root.  On any
finite graph with one root seed and $0<\alpha<1$,
\begin{equation}
 R_\star
 \leq
 1+\frac{
  \log_+\!\bigl((1-\alpha)/(\alpha\rho)\bigr)
 }{
  \log\!\bigl((1+\sqrt\alpha)/(1-\sqrt\alpha)\bigr)
 }
 \leq
 1+\frac{
  \log_+\!\bigl((1-\alpha)/(\alpha\rho)\bigr)
 }{2\sqrt\alpha}.
 \label{eq:graph-support-radius}
\end{equation}
\end{lemma}

\begin{proof}
The spectrum of $Q$ lies in $[\alpha,1]$.  The Chebyshev approximation to
$q\mapsto q^{-1}$ gives, for every $k\geq0$, a degree-$k$ polynomial $p_k$
such that
\begin{equation}
 \norm{Q^{-1}-p_k(Q)}_2
 \leq\frac{2}{\alpha}\lambda_\alpha^{k+1}.
 \label{eq:chebyshev-inverse-error}
\end{equation}
Indeed, with
$c=(1+\alpha)/(1-\alpha)$ and
$q_{k+1}(x)=
T_{k+1}((1+\alpha-2x)/(1-\alpha))/T_{k+1}(c)$, the polynomial
$p_k(x)=(1-q_{k+1}(x))/x$ has degree $k$ and the displayed estimate follows
from
$T_{k+1}(c)=(\lambda_\alpha^{-(k+1)}+
\lambda_\alpha^{k+1})/2$.
If $\operatorname{dist}(o,v)=\ell$ and $k=\ell-1$, sparsity gives
$(p_k(Q))_{vo}=0$.  Therefore, for the unregularized solution
$\bar x:=Q^{-1}b$,
\begin{equation}
 0\leq x_v^\star\leq\bar x_v
 \leq 2\lambda_\alpha^\ell,
 \label{eq:graph-green-decay}
\end{equation}
where obstacle comparison gives the middle inequality; the harmless same
bound at $\ell=0$ follows directly from
$\norm{Q^{-1}}_2\norm{b}_2\leq1$.  Here the single-seed normalization
$b_o=\alpha/\sqrt{d_o}$ was used.

Suppose $R_\star\geq1$ and choose a deepest active vertex $v$.  Its active
KKT equality and nonnegativity give
\begin{equation}
 \alpha\rho\sqrt{d_v}
 <\sum_{u\in\mathcal N(v)\cap S^\star}
   \frac{1-\alpha}{2\sqrt{d_ud_v}}x_u^\star.
 \label{eq:deepest-active-demand}
\end{equation}
Every neighbor of $v$ has distance at least $R_\star-1$.  Applying
\cref{eq:graph-green-decay} and
$\sum_{u\in\mathcal N(v)}d_u^{-1/2}\leq d_v$ bounds the right-hand side by
$(1-\alpha)\sqrt{d_v}\lambda_\alpha^{R_\star-1}$.  Hence
\begin{equation*}
 \lambda_\alpha^{R_\star-1}
 >\frac{\alpha\rho}{1-\alpha}.
\end{equation*}
Taking logarithms proves the first inequality in
\cref{eq:graph-support-radius}.  The second uses
$\log((1+t)/(1-t))\geq2t$ for $0<t<1$.
\end{proof}

\begin{theorem}[Condition-free lazy-response solver on an arbitrary tree]
\label{thm:lazy-tree-response}
Let $\Delta$ be the maximum degree.  There is an exact adjacency-local RPPR
algorithm for every finite tree with one seed.  It discovers
$S^\star(\rho)$ without a support, radius, or confinement oracle and uses
\begin{align}
 \Work_{\rm lazy\text{-}tree}
 &\leq
 O\!\left(
  \vol(S^\star)
  +\log(1+\Delta)
    \sum_{v\in S^\star}\operatorname{dist}(o,v)
 \right)\notag\\
 &\leq
 O\!\left(
  \vol(S^\star)\bigl[1+R_\star\log(1+\Delta)\bigr]
 \right)\notag\\
 &=\widetilde O\!\left(\frac1{\rho\sqrt\alpha}\right)
 \label{eq:lazy-tree-response-work}
\end{align}
degree work and $O(\vol(S^\star))$ storage.  The final equality hides only
the displayed logarithms in $\Delta$ and $(\alpha\rho)^{-1}$.
\end{theorem}

\begin{proof}
Turn \cref{cor:inverse-tree-response} into a lazy iterator.  The first event
of the response from a child $u$ is its activation threshold $\tau_u$, which
requires only the child edge and degree.  Once $u$ activates, scan its
adjacency list once and initialize a heap with the first event of each child.
On the current affine piece, the next breakpoint of $H_u$ is the minimum next
child breakpoint in the common argument $z_u$.  Map that point through the
current affine formula for $H_u$, emit the resulting parent-coordinate event,
update the aggregate slope and intercept, and advance only the child streams
attaining the minimum.  The positivity in
\cref{eq:inverse-response-slope} preserves event order.  Induction from the
leaves proves that this iterator emits exactly the labeled activation
breakpoints of $\mu_u$, including tied labels, while never inspecting a
descendant of an inactive boundary vertex.

At the root, merge the child iterators in the same fashion and walk the
piecewise-affine left side of \cref{eq:root-response-equation}.  Its slope is
positive by \cref{eq:root-response-slope-bound}.  On each piece, compute its
affine zero.  If that zero occurs no later than the next event, stop; otherwise
consume the event and continue.  A label is consumed exactly when its vertex
is positive at the final root value.  Thus the consumed labels, together with
the active root, are exactly $S^\star(\rho)$.  Equality with the next event is
not consumed, consistently with the definition of support.  If
$c_{\rho,o}\leq0$, the zero solution is certified before any event.

Building the heaps scans each active adjacency list once, for total
$O(\vol(S^\star))$ work and storage, including inactive boundary edges.  An
event labeled by $v$ is advanced once through the heap of each ancestor of
$v$, costing
$O(\operatorname{dist}(o,v)\log(1+\Delta))$.  Iterator state stores only the current affine
piece and one next event per exposed child, so response histories are not
copied through all ancestors.  After discovery, apply
\cref{thm:forest-resolution} once to $S^\star$ and check its boundary KKT
certificate in $O(\vol(S^\star))$ work.  Summing the event charges proves the
first two lines of \cref{eq:lazy-tree-response-work}; the support-volume bound
and \cref{lem:graph-support-radius} prove the last.
\end{proof}

This closes the nonsymmetric-tree free boundary at the intended product
scale.  It does not claim a universal lower bound: paths and symmetric
spiders remain output-linear special cases.  Cycles destroy vertex-by-vertex
subtree independence, but not the underlying one-dimensional event order.
